\chapter{Planificación}
\label{chap:planificacion}

\section{Alcance del proyecto}
\subsection{Definición del alcance}
Una aplicación de gestión de proyectos profesional a día de hoy es una herramienta con muchísimas funcionalidades, y es imposible abarcar todo el espectro de posibilidades en un proyecto de estas características, tanto en fecha como en dedicación. 

Por ello, es necesario definir un alcance claro y realista para el proyecto, que permita centrarse en las funcionalidades indispensables en aplicaciones de este este tipo, añadiendo la innovación planteada en el proyecto, y dejando de lado funcionalidades que, aunque interesantes, no son imprescindibles para el correcto funcionamiento de la aplicación.

Para definir el alcance del proyecto y priorizar las funcionalidades desde el punto de vista del usuario final, se ha utilizado la técnica de las Historias de Usuario (User Stories), que permite describir las funcionalidades desde la perspectiva del usuario, y priorizarlas en función de su importancia para el usuario final.
Para ello, el primer paso es identificar a los usuarios que van a utilizar la aplicación.


\subsection{Exclusiones}
Tras recopilar las historias de usuario, se han identificado algunas funcionalidades que, aunque podrían ser interesantes, quedan fuera del alcance del proyecto por no ser imprescindibles para el correcto funcionamiento de la aplicación, o por requerir un esfuerzo de desarrollo desproporcionado en relación al beneficio que aportan. De esta forma, se han desarrollado las historias de usuario con una prioridad de \emph{Must Have} y \emph{Should Have}, dejando fuera funcionalidades con prioridad de \emph{Could Have} o \emph{Won't Have}.

Además, por lo referente a la autenticación, se ha decidido limitar el proyecto a la autenticación mediante OAuth con Google, dejando fuera otras opciones de autenticación (Microsoft, Facebook, X, etc.) ya que el foco del proyecto no es la integración de múltiples proveedores de autenticación.

También se han dejado fuera funcionalidades relacionadas con la gestión de recursos (gestión de archivos, integración con repositorios de código, etc.) que aunque serían muy útiles e interesantes, no pueden ser implementadas dentro del alcance del proyecto.

\section{Cronograma}
